\section{练习}

\begin{question}\begin{drill}[电子流轰击接地金属球]\label{e2:drill:1}
半径 $a$ 的孤立金属球经电阻 $R$ 接地，电子从很远处以速率 $v$ 落到球上，稳态每秒实际到达 $N$ 个。求球电荷、电阻发热功率和球吸收碰撞能的功率，忽略其他能量交换。
\end{drill}
\typeout{ECFIG:fig:e2:drill:1:s}\fig{fig:e2:drill:1:s}{电子实际落到半径 $a$ 的球上；接地电阻为 $R$，两者符号分开}{\begin{tikzpicture}

\draw(0,0)circle(1.5);\fill(0,0)circle(1.5pt)node[below]{$O$};\draw[<->](0,0)--(-1.06,1.06)node[midway,left]{$a$};
\draw(0,-1.5)to[R,l={$R$}](0,-3.8)node[ground]{};
\foreach \y in {-.5,0,.5}{\draw[flow](-4,\y)--(-2,\y);}\node at(-3,1.2){$N$ 个/秒，速率 $v$};

\end{tikzpicture}}
\end{question}
\answerlink{e2:drill:1}
\begin{question}\begin{drill}[三层同心球体依次连接]\label{e2:drill:2}
半径依次为 $a<b<c$ 的导体球 $A$、球壳 $B,C$ 同心。初始 $A$ 电势 $V_A$，$B$ 净电荷零，$C$ 净电荷 $Q_C$。先接通 $A,B$，再接通 $B,C$。求两次 $A$ 电势增量。
\end{drill}
\typeout{ECFIG:fig:e2:drill:2:s}\fig{fig:e2:drill:2:s}{题设：三层同心球体依次连接}{\begin{tikzpicture}[scale=.88]
\draw[world](0,0)circle(.55);\draw[world](0,0)circle(1.15);\draw[world](0,0)circle(1.75);\node at(0,0){$A$};\node at(.8,.5){$B$};\node at(1.5,.8){$C$};\node at(3.2,0){$a<b<c$};
\end{tikzpicture}}
\end{question}
\answerlink{e2:drill:2}
\begin{question}\begin{drill}[电势为零与电场为零]\label{e2:drill:3}
等量电荷 $+q,-q$ 固定于 $x=-a,+a$。求原点电势和电场。若改为两个 $+q$，再求这两个量，并解释结论。
\end{drill}
\typeout{ECFIG:fig:e2:drill:3:s}\fig{fig:e2:drill:3:s}{题设：电势为零与电场为零}{\begin{tikzpicture}[scale=.88]
\draw[world,fill=white](-1.5,0)circle(.12);\node[above]at(-1.5,0.15){$+q$};\draw[world,fill=white](1.5,0)circle(.12);\node[above]at(1.5,0.15){$-q$};\draw[axis](-3,0)--(3,0)node[right]{$x$};\node[below]at(0,0){$O$};
\end{tikzpicture}}
\end{question}
\answerlink{e2:drill:3}
\begin{question}\begin{drill}[三点电荷的装配与释放]\label{e2:drill:4}
边长 $a$ 的等边三角形三顶点各带正电荷 $q$。先固定其中两电荷，将第三个质量 $m$ 的电荷从静止释放。求初始系统相互作用能，以及该电荷到无穷远的速率。
\end{drill}
\typeout{ECFIG:fig:e2:drill:4:s}\fig{fig:e2:drill:4:s}{题设：三点电荷的装配与释放}{\begin{tikzpicture}[scale=.88]
\draw[world](0,0)--(4,0)--(2,3.464)--cycle;\node[below]at(0,0){$q$};\node[below]at(4,0){$q$};\node[above]at(2,3.464){$q,m$};\node at(2,-.7){每边 $a$};
\end{tikzpicture}}
\end{question}
\answerlink{e2:drill:4}
\begin{question}\begin{drill}[圆盘轴线上的电势与场跳变]\label{e2:drill:5}
半径 $R$ 的圆盘均匀带面密度 $\sigma>0$。以无穷远为零，求轴线全体 $x$ 的电势；用导数核验两侧电场，并说明电势是否跳变。
\end{drill}
\typeout{ECFIG:fig:e2:drill:5:s}\fig{fig:e2:drill:5:s}{题设：圆盘轴线上的电势与场跳变}{\begin{tikzpicture}[scale=.88]
\draw[world](0,0)ellipse(.5 and 1.4);\draw[axis](-1,0)--(4.5,0)node[right]{$x$};\fill(0,0)circle(1.3pt)node[below left]{$O$};\fill(3,0)circle(1.6pt)node[above]{$P$};\draw[dim](0,-1.8)--(3,-1.8)node[midway,below]{$x$};\draw[dim](0,0)--(0,1.4)node[midway,left]{$R$};\node at(-1.3,1.9){均匀圆盘 $\sigma$};
\end{tikzpicture}}
\end{question}
\answerlink{e2:drill:5}
\begin{question}\begin{drill}[有限线段的中垂面电势]\label{e2:drill:6}
长 $2l$ 的线段位于 $y$ 轴，线密度 $\lambda$，求 $P=(x,0),x>0$ 处电势；由求导核对中垂线场，并讨论无限长极限的零点。
\end{drill}
\typeout{ECFIG:fig:e2:drill:6:s}\fig{fig:e2:drill:6:s}{题设：有限线段的中垂面电势}{\begin{tikzpicture}[scale=.88]
\draw[axis](-.2,0)--(4,0)node[right]{$x$};\draw[axis](0,-.2)--(0,2.4)node[above]{$y$};\draw[world,very thick](0,-2)--(0,2);\fill(3,0)circle(1.5pt)node[above]{$P$};\node[left]at(0,2){$l$};\node[left]at(0,-2){$-l$};
\end{tikzpicture}}
\end{question}
\answerlink{e2:drill:6}
\begin{question}\begin{drill}[电偶极子的远场电势]\label{e2:drill:7}
等量异号电荷间距 $2a$，偶极矩 $\bm p$ 从负指向正。场点距中点 $r\gg a$，与偶极轴夹角 $\theta$。求电势首阶远场，并由球坐标方向导数得到轴线、中垂线场。
\end{drill}
\typeout{ECFIG:fig:e2:drill:7:s}\fig{fig:e2:drill:7:s}{题设：电偶极子的远场电势}{\begin{tikzpicture}[scale=.88]
\draw[guide](0,0)--(2,0);
\draw[world,fill=white](-1,0)circle(.12);\node[above]at(-1,.15){$-q$};
\draw[world,fill=white](1,0)circle(.12);\node[below]at(1,-.16){$+q$};
\fill(0,0)circle(1.2pt)node[below left]{$O$};
\draw[flow](-.5,-.6)--(.5,-.6)node[midway,below]{$\bm p$};
\draw[guide](0,0)--(3,2.3)node[midway,above]{$r$};
\draw(.55,0)arc(0:37.475:.55);\node at(.8,.25){$\theta$};
\fill(3,2.3)circle(1.5pt)node[right]{$P$};
\draw[dim](-1,-1.3)--(1,-1.3)node[midway,below]{$2a$};
\end{tikzpicture}}
\end{question}
\answerlink{e2:drill:7}
\begin{question}\begin{drill}[带电球壳分裂后的能量比较]\label{e2:drill:8}
一个绝缘薄球壳半径 $R$、总电荷 $Q$，均匀分布。若外界将其电荷重新均匀分配到两个相距无限远、半径也为 $R$ 的相同球壳上，每壳 $Q/2$，求静电能变化和外界缓慢做功。
\end{drill}
\typeout{ECFIG:fig:e2:drill:8:s}\fig{fig:e2:drill:8:s}{题设：带电球壳分裂后的能量比较}{\begin{tikzpicture}[scale=.88]
\draw[world](0,0)circle(1.5);\fill(0,0)circle(1.2pt)node[below]{$O$};\draw[dim](0,0)--(-1.06,1.06)node[midway,left]{$R$};\node at(2.8,.5){$Q,R$};
\end{tikzpicture}}
\end{question}
\answerlink{e2:drill:8}
\begin{question}\begin{drill}[电子在电势差中的速度]\label{e2:drill:9}
电子质量 $m$、电荷 $-e$，从电势 $V_A$ 处以速率 $v_0$ 运动到 $V_B$，仅受静电力。求末速与可达条件；若初静止，判断电子趋向高势还是低势。
\end{drill}
\typeout{ECFIG:fig:e2:drill:9:s}\fig{fig:e2:drill:9:s}{题设：电子在电势差中的速度}{\begin{tikzpicture}[scale=.88]
\draw[world](0,-1)--(0,1)(4,-1)--(4,1);\node[above]at(0,1){$V_A$};\node[above]at(4,1){$V_B$};\draw[flow](.3,0)--(1.6,0)node[midway,above]{$v_0,-e$};
\end{tikzpicture}}
\end{question}
\answerlink{e2:drill:9}
